% Notebook 21 concrete module. No preamble. Prefix cv-.
\section{Concrete classical vertices uniform through refinement}
\label{cv-section}

The fine spacing is $\eta=M^{-1}$, with $M=L^j$, fixed odd
$L$, dimension four and current spacing one. Fine pairings
use $\eta^4$ times counting measure; current pairings use
counting measure. All constants below are uniform in $M$ and
compatible periodic volume. The conclusion is a set of concrete
inputs for a rooted support expansion, not an assumption that
one weighted response estimate already gives a sum over all
input insertions.

\subsection{Complete physical field spaces and known bounds}
\label{cv-spaces-section}

Fix $0<\alpha<1$ and use complete fine-field spaces
\begin{equation}\label{cv-spaces}
 \begin{split}
 X_\eta=\mathcal N_\alpha:\quad
 \norm A_X&=\norm A_\infty+\norm{\nabla_\eta A}_\infty
                         +[\nabla_\eta A]_\alpha,\\
 Y_\eta&=C^\alpha\text{ of fine jet coefficients},\\
 Z_\eta&=\ell^\infty\text{ of fine coefficients},\qquad
 E=\ell^\infty\text{ of current bond fields}.
 \end{split}
\end{equation}
Local H\"older norms use physical distances and fixed cell
collars. Color and orientation spaces have fixed finite
dimension. The plaquette jet $S=(I,\nabla_\eta)$, with the
needed finite component selections, maps $X_\eta$ to $Y_\eta$
uniformly. No sharp cell indicator is differentiated.

The proved Wilson-force module gives
\[
 \mathcal S_\eta(A)=\eta^4\sum_p\mathcal L_\eta((SA)_p),
 \qquad N(A)=S^Tf_\eta(SA),
\]
where $\mathcal L_\eta$ is uniformly analytic in the jet,
its quadratic part is the Maxwell density, and $f_\eta$ starts
at degree two. The full corrected constraint of Notebook13 is
\[
 \mathcal F(A)=QA+\mathcal R(A),\qquad Q=Q_M^{\rm st}.
\]
It is analytic on a refinement-uniform supremum ball and has
a fixed physical dependency radius. Notebook15 also proves
uniform holomorphy of $D\mathcal R(A)^T$ in the physical
bounded-row operator norm, with norm $O(\norm A_\infty)$.
The transpose includes the physical fine measure.

The weighted complete-field bounds for $H,\mathcal C$ and
their differentiated compositions are recorded explicitly in
Section~\ref{cv-linear-section}. They include the
equal-exponent constrained Schauder theorem of Notebook15.
The scalar action of the ordinary homogeneous minimizer is
exactly the action at the corrected fixed-slice minimizer by
Notebook15's gauge transfer. This does not identify the two
constraint derivatives.

The underlying corrected one-step coordinate map is specified
at Raw I, lines 24320--24325 and 24344--24369, with the
ordinary/corrected bridge at 47076--47087. Its physical
scaling is $(Lh)^{-1}F_L(hA)$. The Wilson density and pairing
use the conventions already matched in Notebooks14--15.
No additional claim about the full original orbit domains is
made here.

\subsection{Smooth cell spaces without an extension assumption}
\label{cv-cells-section}

Choose a smooth partition of unity $\{\chi_i\}$ indexed by
unit current cells, with support in fixed enlarged cells
$S_i$, uniformly bounded derivatives through order three, and
uniformly finite overlap. For example, divide a nonnegative
compactly supported smooth bump by the sum of its integer
translates. Periodize on compatible tori, counting repeated
supports with multiplicity.

For each fine field type let $E_i$ be the Banach subspace of
\emph{complete} fields supported in $S_i$, with its global
$X_\eta$, $Y_\eta$ or $Z_\eta$ norm. Write $I_i$ for inclusion
and $P_iu=\chi_i u$. For current fields use coordinate
restriction and insertion in one current cell. Then
\begin{equation}\label{cv-partition}
 u=\sum_iI_iP_iu,\qquad
 \norm{P_iu}_{E_i}\le C\norm u_{\text{fixed collar of }i}.
\end{equation}
The reconstruction
\[
 \mathscr R((u_i))=\sum_iI_iu_i
\]
is bounded from the block supremum space to the complete
global space, uniformly in $\eta$ and volume. In particular,
$\mathscr RP=I$ and $P\mathscr R$ is a bounded projection.
Arbitrary restrictions on a cell are not extended by an
unproved extension theorem; the spaces $E_i$ already consist
of supported complete fields.

For the only delicate norm $X_\eta$, the exact product rule is
\begin{equation}\label{cv-product-rule}
 \nabla_\eta(\chi A)(x)
 =\chi(x+\eta e)\nabla_\eta A(x)
          +(\nabla_\eta\chi)(x)A(x).
\end{equation}
The sampled smooth functions and their first differences have
uniform physical $C^\alpha$ norms. Telescoping controls the
local $C^\alpha$ norm of $A$ by $\norm A_X$ without an inverse
mesh factor. The H\"older product inequality proves
\eqref{cv-partition}. Finite overlap, also for pairs of points
in fixed collars, proves the reconstruction bound. The
arguments for $Y_\eta$ and $Z_\eta$ are simpler.

\subsection{A summed mixed-insertion vertex bound}
\label{cv-vertex-section}

For a finite set of current cell labels $S$, let $\tau(S)$ be
the minimum length of a connected current-lattice graph
containing $S$, allowing unmarked auxiliary vertices. This is
graph Steiner length. It is not the minimum-spanning-tree
length on the terminals alone.

For a local analytic map $F$, expand after reconstruction:
\begin{equation}\label{cv-vertex-expansion}
 P_iF(\mathscr Ru)=
 \sum_{n\ge0}\sum_{j_1,\ldots,j_n}
 V_{n;i;j_1,\ldots,j_n}[u_{j_1},\ldots,u_{j_n}].
\end{equation}
The coefficients contain $1/n!$. A fixed finite direct sum
handles several field types. Define
\begin{equation}\label{cv-vertex-norm}
 v_n(\kappa)=\sup_i\sum_{j_1,\ldots,j_n}
 e^{\kappa\tau(\{i,j_1,\ldots,j_n\})}
             \norm{V_{n;i;j_1,\ldots,j_n}}.
\end{equation}
The coefficient norm is the multilinear operator norm of
the cell Banach spaces. All input cell labels are summed.

\begin{lemma}[Finite-range analytic vertices]\label{cv-local-lemma}
Suppose each output reads only labels within distance $R$,
at most $D_R$ labels, and $F$ is analytic and uniformly
bounded by $M_F$ on a global ball of radius $r_F$. If
$\norm{\mathscr R}\le C_P$, then
\begin{equation}\label{cv-local-majorant}
 v_n(\kappa)\le M'_F
       \left(\frac{eC_PD_Re^{\kappa R}}{r_F}\right)^n,
\end{equation}
with a fixed output-localization factor in $M'_F$.
If $F$ starts at degree two, the generating majorant is
analytic and quadratic at zero.
\end{lemma}

\begin{proof}
For normalized directions, multivariable Cauchy with radii
$r_F/(nC_P)$ bounds the derivative by
$M'_F(nC_P/r_F)^n$. Division by $n!$ and
$n^n/n!\le e^n$ gives the coefficient bound. There are at
most $D_R^n$ nonzero ordered tuples. Joining the root to
each input by a shortest path costs at most $nR$, proving
\eqref{cv-local-majorant}. This argument never counts
individual fine coordinates. On a smaller analytic margin
it also works at a nonzero fixed background, though the
constant and linear terms then need not vanish.
\end{proof}

Rooting at an input and summing the output and other inputs
has the same bound up to fixed factors, since the possible
outputs lie in a fixed neighborhood of that input.
For support-family substitution the exact geometric bound is
\begin{equation}\label{cv-tree-union}
 \tau\left(\{i\}\cup\bigcup_rX_r\right)
 \le\tau(\{i,j_1,\ldots,j_n\})
                    +\sum_r\tau(\{j_r\}\cup X_r).
\end{equation}
Join the graphs at their roots. The $j_r$ remain allowed as
unmarked auxiliary vertices when they are no longer terminals.
This constant-one inequality can fail for metric
minimum-spanning-tree length. If that convention is denoted
by $\ell_{\rm MST}$, a doubled traversal and shortcut give
\begin{equation}\label{cv-mst-conversion}
 \ell_{\rm MST}(S)\le2\tau(S).
\end{equation}
Thus a final weight $e^{\mu\ell_{\rm MST}}$ is covered by
a Steiner exponent at least $2\mu$. That spatial reserve
must be retained at the polymer interface.

\paragraph{Arbitrary fixed cardinality reserve.}
The concrete coefficients here are expanded at the fixed
flat background. They introduce no extra variable support.
For any prescribed $a\ge0$,
\begin{equation}\label{cv-cardinality}
 e^{a|\cup_rX_r|}\le\prod_re^{a|X_r|}.
\end{equation}
Consequently adding $e^{a|X|}$ to the rooted support norm
does not change the nonlinear vertex majorants. Linear
operators change the root, not the variable support. Only
the initial source $B_j$, with support $\{j\}$, costs
$e^a$. The source radius may shrink with $a$ while remaining
uniform in refinement and volume. This can pay a fixed
support-cardinality factor in a later polymer norm; small
spatial decay alone would not pay that factor. Hiding a
background-dependent minimizer in the Taylor coefficients
would require a new support argument and is not done here.

\subsection{The actual nonlinear vertices}
\label{cv-concrete-section}

\paragraph{The Wilson jet coefficient.}
For a positive plaquette use
$z=(a,b,c,e)=(A_\mu,A_\nu,\nabla_\mu A_\nu,
\nabla_\nu A_\mu)$ at its base point. The exact expressions are
\begin{equation}\label{cv-wilson-jet}
 \begin{split}
 P_\eta(z)&=e^{i\eta a}e^{i\eta b+i\eta^2c}
                 e^{-i\eta a-i\eta^2e}e^{-i\eta b},\\
 \mathcal L_\eta(z)&=-\frac12\tau\left[
       \left(\frac{P_\eta(z)-I}{\eta^2}\right)^2
                                      P_\eta(z)^{-1}\right].
 \end{split}
\end{equation}
The singularity at $\eta=0$ is removed before differentiating
vertices: $P_0=I$ and $\partial_\eta P_\eta|_0=0$.
Joint analyticity on a fixed jet ball gives uniform
coefficients. The bounded jet map and the $C^\alpha$
product estimate make
\[
 F_\eta^{\rm nl}(A)=f_\eta(SA):X_\eta\longrightarrow Y_\eta
\]
analytic on a uniform ball and quadratic at zero. Each
plaquette output reads a bounded physical neighborhood.
Lemma~\ref{cv-local-lemma} therefore proves its full
mixed-insertion bound uniformly in $\eta$.
The transpose $S^T$ stays outside this vertex. A uniform
bounded-field estimate for the unpropagated fine force
$S^TF_\eta^{\rm nl}$ is neither used nor claimed.

\paragraph{The Wilson energy density.}
Set
\begin{equation}\label{cv-cell-energy}
 \mathcal E_{i,\eta}(A)=\eta^4\sum_{p>0}
       \chi_i(x_p)\mathcal L_\eta((SA)_p).
\end{equation}
Then $\mathcal S_\eta(A)=\sum_i\mathcal E_{i,\eta}(A)$
exactly. Although a current cell contains $O(\eta^{-4})$
fine plaquettes,
\begin{equation}\label{cv-cell-mass}
 \eta^4\sum_{p>0}|\chi_i(x_p)|\le C_d.
\end{equation}
The fine differences are already bounded in the jet norm.
Thus every cell energy is a uniformly bounded analytic
scalar map on a fixed $\mathcal N_\alpha$ ball, with fixed
current dependency range. Its coefficients satisfy
\eqref{cv-local-majorant}. Replacing $\mathcal L_\eta$ by
its remainder after the quadratic Maxwell density gives
the same conclusion starting at degree three. Dropping
$\eta^4$, or treating all fine differences as unrelated
bounded coordinates, would lose this estimate.
The ordinary constrained scalar action is the same sum at
the corrected minimizer by the established gauge transfer.

\paragraph{The full corrected constraint.}
At input spacing $h$ and output spacing $Lh$,
\begin{equation}\label{cv-one-step}
 \begin{split}
 \mathcal F_h(A)&=(Lh)^{-1}F_L(hA)=Q_hA+\mathcal R_h(A),\\
 \norm{\mathcal R_h(A)}_\infty&\le Ch\norm A_\infty^2,
 \qquad
 \norm{D\mathcal R_h(A)}_{\infty\to\infty}
       \le Dh\norm A_\infty.
 \end{split}
\end{equation}
For $h_t=L^{t-j}$, $\sum_{t<j}h_t\le(L-1)^{-1}$.
The proved Riccati bound keeps every intermediate physical
field in a common ball. Hence the composed remainder
$\mathcal R$ is uniformly analytic, quadratic, and has
physical dependency radius at most
$r_L\sum_th_t\le R_*$.
It follows that $\mathcal R:X_\eta\to E$ satisfies
Lemma~\ref{cv-local-lemma}. The scale sum is already included
in this uniform bound; higher coefficients are obtained by
Cauchy in that ball, not by an uncontrolled count of
multilevel diagrams.

\paragraph{The physical-adjoint constraint vertex.}
Define
\begin{equation}\label{cv-adjoint-vertex}
 \mathcal V_\eta(A,m)=D\mathcal R(A)^Tm:
                X_\eta\times E\longrightarrow Z_\eta.
\end{equation}
Ordinary holomorphy of $\mathcal R$ alone would not control
this physical transpose. The proved row estimate does:
at one step,
\[
 \norm{D\mathcal R_h(A)^T}_{\infty\to\infty}
       \le D_*h\norm A_\infty,\qquad
 D_*=N_LK_0L^{d-1}.
\]
The factor $L^d$ in the physical transpose is included here.
Both $Q_h$ and $Q_h^T$ contract the bounded norm.
Transposing the composition and using the intermediate
field bound proves
\begin{equation}\label{cv-adjoint-bound}
 \norm{D\mathcal R(A)^T}_{\infty\to\infty}
    \le \exp\left(\frac{2D_*\norm A_\infty}{L-1}\right)-1.
\end{equation}
The map is holomorphic into this physical absolute-row
operator space. Thus \eqref{cv-adjoint-vertex} is jointly
analytic and $O(\norm A_\infty\norm m_\infty)$.
Its output at a fine bond reads coarse $m$ within $R_*$,
and background $A$ within $2R_*$, because each relevant
coarse equation contains that bond and reads its own
neighborhood. Lemma~\ref{cv-local-lemma} applies with
total degree starting at two. Cauchy is taken in the
physical-row operator norm; it is not applied entrywise
and followed by a count of fine coefficients.
The output is only required to be bounded: it enters
$\mathcal C$ or $H^T$, and no derivative of $m$ is taken.

\subsection{Summable linear blocks, including the derivative kernels}
\label{cv-linear-section}

The established complete-field estimates, for every positive
weight with $w(i)/w(j)\le e^{\kappa_0d(i,j)}$, are
\begin{equation}\label{cv-linear-inputs}
 \begin{gathered}
 H:E_w\to X_{\eta,w},\qquad
 \mathcal C:Z_{\eta,w}\to X_{\eta,w},\qquad
 \mathcal CS^T:Y_{\eta,w}\to X_{\eta,w},\\
 H^T:Z_{\eta,w}\to E_w,\qquad
 (SH)^T:Z_{\eta,w}\to E_w,\qquad q:E_w\to E_w.
 \end{gathered}
\end{equation}
Restricting the fifth input to $Y_\eta$ is sufficient.
The third estimate has the same $\alpha$ at input and
output; a positive H\"older loss would not suffice.
These are the constrained Schauder and extension bounds
already used in the weighted response proof. They do not
require the gauge projection $\Pi$ to act on arbitrary
bounded fields.

For any such $T$, form $T_{ij}=P_iTI_j$. Given a supported
input $u_j$, use $w_j(n)=e^{\kappa_0d(n,j)}$. Its weighted
norm is bounded by its unweighted norm times a fixed collar
factor. The complete weighted estimate gives
\begin{equation}\label{cv-block-envelope}
 \norm{T_{ij}}\le C_Te^{-\kappa_0d(i,j)}.
\end{equation}
This is an operator norm between cell Banach spaces, not an
absolute integral estimate for a singular differentiated
fine-coordinate kernel. In particular its diagonal block is
controlled by the equal-exponent Schauder bound.
For every $0<\kappa<\kappa_0$, lattice growth then gives
\begin{equation}\label{cv-block-schur}
 \max\left\{
  \sup_i\sum_je^{\kappa d(i,j)}\norm{T_{ij}},\,
  \sup_j\sum_ie^{\kappa d(i,j)}\norm{T_{ij}}\right\}
                \le C_{T,\kappa}.
\end{equation}
Both sums follow from the common envelope, without any
self-adjointness assumption in these Banach norms.
Periodic-distance sums have the same uniform upper bound.

The decomposition is exact; for example,
\begin{equation}\label{cv-exact-linear-splitting}
 \mathcal CS^Tf=\sum_j\mathcal CS^T(\chi_jf).
\end{equation}
Every $\chi_jf$ is a complete $C^\alpha$ field. Smooth
cutoff derivatives are included in $S^T(\chi_jf)$ and
controlled by the same estimate; they cancel in the sum.
No sharp splitting is used. The maps $H^T$ and $(SH)^T$
likewise retain their physical transpose normalization.
Output first differences and their H\"older seminorms are
already in $X_\eta$, while source derivatives occur only
through $\mathcal CS^T$ and $(SH)^T$. Separate absolute
fine-coordinate derivative bounds are unnecessary for
\eqref{cv-vertex-norm}.

\subsection{An exact simultaneous fixed point using these vertices}
\label{cv-fixed-point-section}

The corrected constraint is $\mathcal F(A)=QA+\mathcal R(A)$.
Let $m$ be the multiplier of Notebook15 and write
$\mathcal V_\eta(A,m)=D\mathcal R(A)^Tm$.
The exact stationarity equations are equivalent near zero to
\begin{equation}\label{cv-fixed-point}
 \begin{split}
 A={}&HB-H\mathcal R(A)
          -\mathcal CS^Tf_\eta(SA)+\mathcal C\mathcal V_\eta(A,m),\\
 m={}&qB-q\mathcal R(A)
          +(SH)^Tf_\eta(SA)-H^T\mathcal V_\eta(A,m).
 \end{split}
\end{equation}
Indeed the second equation rearranges to
\[
 (I+H^TD\mathcal R(A)^T)m
        =q[B-\mathcal R(A)]+H^TN(A),
\]
and the first is $A=H[B-\mathcal R(A)]+\xi$ with
$\xi=-\mathcal CN(A)+\mathcal CD\mathcal R(A)^Tm$.
These are the proved constraint chart and reduced
stationarity identities. Conversely the exact constraints
of $\mathcal C$ put $\xi$ in the fixed flat slice, and the
small multiplier inverse recovers the earlier formulation.
No extra small critical points are introduced.

All nonlinear primitives in \eqref{cv-fixed-point} start at
degree two: $\mathcal R(A)$, $f_\eta(SA)$, and
$D\mathcal R(A)^Tm$. Their mixed-insertion majorants are
analytic and uniform by the preceding proof, while their
linear propagators satisfy \eqref{cv-block-schur}.
Apply $P$ to the equations and reconstruct with $\mathscr R$.
A fixed point in compatible cell components reconstructs
exactly to \eqref{cv-fixed-point}; the overlap constants are
absorbed by reducing the small radius. This supplies the
concrete inputs for the separate rooted support-family
fixed-point theorem, including any fixed cardinality reserve.
Tree summation, canonical grouping of supports and the final
scalar-action expansion remain conclusions of that theorem,
not substitutes for its proof.

\subsection{An obstruction to using the full ordinary adjoint}
\label{cv-ordinary-section}

The full ordinary rooted average is analytic from small
complete physical connections to bounded current data, with
fixed physical dependency range, by the ordinary trajectory
theorem. It therefore has $X_\eta\to E$ vertex bounds.
Its physical adjoint on arbitrary bounded current fields is
different and is not uniformly bounded through refinement.

Let $Q_{\rm ord}$ be its flat derivative in physical
potentials. Exact ordinary gauge covariance gives
\begin{equation}\label{cv-ordinary-gradient}
 Q_{\rm ord}D_\eta\lambda
          =D_1(\lambda|_{\text{final centers}}),
\end{equation}
with the same sign on both sides for either consistent
gauge convention. On a sufficiently large torus, take a
unit color vector $\lambda$ at one final-center fine site
and zero at every other fine site. Its fine gradient is
nonzero on $2d$ bonds, each of magnitude $\eta^{-1}$, so
\begin{equation}\label{cv-point-gradient-norms}
 \norm{D_\eta\lambda}_{2,\eta}
       =\sqrt{2d}\,\eta^{d/2-1},\qquad
 \norm{D_\eta\lambda}_{1,\eta}=2d\eta^{d-1}.
\end{equation}
The current gradient in \eqref{cv-ordinary-gradient} has
magnitude one on $2d$ bonds. Therefore
\begin{equation}\label{cv-ordinary-hilbert-obstruction}
 \norm{Q_{\rm ord}}_{L^2_\eta\to L^2_1}
                     \ge\eta^{1-d/2}.
\end{equation}
Choose $m$ to equal that current gradient, so
$\norm m_\infty=1$ and
$\langle Q_{\rm ord}D_\eta\lambda,m\rangle_1=2d$.
Physical transposition and
\eqref{cv-point-gradient-norms} imply
\begin{equation}\label{cv-ordinary-adjoint-obstruction}
 \norm{Q_{\rm ord}^T}_{\infty\to\infty}\ge\eta^{1-d}.
\end{equation}
In dimension four these are $\eta^{-1}$ and $\eta^{-3}$.
The complete $\mathcal N_\alpha$ norm of the test gradient
is large, so there is no contradiction with
$Q_{\rm ord}:X_\eta\to E$ boundedness.

The corrected flat derivative instead satisfies
$Q_M^{\rm st}D_\eta=D_1Q_0$ with the scalar block mean,
and its physical transpose contracts the bounded norm.
This is why \eqref{cv-fixed-point} uses the corrected map.
Possible cancellations in a specially projected ordinary
adjoint are not assumed. The ordinary current-step pivot
chart of Notebook19 is unaffected: that ratio is the fixed
$L$, so its constants may depend on $L$. The obstruction
concerns unbounded total ratio $M$.

\subsection{Scope of the concrete input}
\label{cv-scope}

The proved inputs are uniformly analytic finite-current-range
Wilson jet and energy vertices, corrected-constraint and
physical-adjoint constraint vertices, summed mixed-insertion
norms, and exponentially summable linear cell blocks including
the derivatives used in the simultaneous equations. The
one-weight theorem is used only to form linear block
envelopes; the nonlinear coefficient proof really sums all
input labels.

No raw fine-coordinate absolute Taylor norm without physical
normalization or regularity is asserted. No ordinary rooted
adjoint bound, independent-current source extension,
mixed-history theorem or quantum integration result is
inferred. Backgrounds remain in the homogeneous small global
potential branch. The scalar gauge transfer, rather than a
false identification of ordinary and corrected derivatives,
matches its classical action to the ordinary branch.
